% PDF 186--191；原书第 157--162 页。
\section{计算高效微调}
\label{sec:ch04-compute-efficient}
PEFT 只训练模型中的一小部分参数，以降低微调大语言模型的开销。计算高效微调（compute-efficient fine-tuning）着眼于量化方法，减少在给定参数数量下微调或推理所需的内存。这些方法通常能在训练和推理开销与任务表现之间取得更好的折中。与 LoRA 等常用方法相比，有些方法会损失一些任务表现，另一些则既提高计算资源的利用效率，又改善任务表现，达到或接近当时的最佳水平。

\subsection{大语言模型量化}
\label{sec:ch04-quantization}
量化（quantization）是一种模型压缩技术，用信息位数较少的形式表示参数，从而缩小模型的总体积（\refcite{ch04-zhao}{Zhao et al., 2023}）。这会减少推理所需的计算资源。量化通常施加在 Transformer 注意力层和前馈层的参数权重上，因为这些层的矩阵乘法运算占大语言模型推理内存消耗的 95\% 以上。改变其中使用的数据类型，可以大幅减少内存消耗（\refcite{ch04-dettmers-2022}{Dettmers et al., 2022}）。

\begin{table}[!htbp]
\centering
\caption{大语言模型中的常用数据类型。表中区分源自其他计算领域的标准数据类型和为机器学习优化的表示，并列出其他常用名称和存储所需位数。}
\label{tab:ch04-types}
\small
\begin{tabularx}{\linewidth}{@{}l l X r@{}}
\toprule
数据类型 & 类型来源 & 其他名称 & 位数\\
\midrule
float32 & 标准 & FP32，单精度浮点格式 & 32\\
float16 & 标准 & FP16，半精度浮点格式 & 16\\
bfloat16 & 机器学习优化 & BF16，Brain 浮点格式 & 16\\
INT8 & 标准 & — & 8\\
INT4 & 标准 & — & 4\\
NF4 & 机器学习优化 & — & 4\\
\bottomrule
\end{tabularx}
\end{table}

数据精度与计算效率之间需要权衡。表~\ref{tab:ch04-types} 列出神经网络，尤其是大语言模型常用的几种数据类型的位宽。例如，将参数权重从 float32 量化为 int8，可以将所需内存降至接近四分之一，实际幅度取决于哪些层和参数不参与量化。这个差别很重要，因为一些模型推理时需要的工作内存，远超多数当时先进 GPU 所能提供的显存。例如，具有 1750 亿参数的 GPT-3 在 float16 精度下需要 325~GB 的存储空间，只能在复杂的多 GPU 集群上运行，无法在更常见的硬件上运行，例如单张配备 80~GB 显存的 NVIDIA A100（\refcite{ch04-frantar}{Frantar et al., 2023}）。

大语言模型的量化大体有两种方案，即训练后量化（post-training quantization，PTQ）与量化感知训练（quantization-aware training，QAT）。先考察大语言模型上几项有影响的训练后量化应用，重点是能够把推理的资源消耗降到 NVIDIA A100 或 NVIDIA A600 等常见硬件可承受的范围，并能取得接近同一模型未量化版本的任务表现的方法。

随后讨论微调中的量化感知训练。这时，对微调适配器做量化感知训练，预训练大语言模型本身不参与这种训练。这类应用也提高了大语言模型推理的计算资源利用效率，使预算有限或受到其他资源约束的实践者能够使用它们，例如在微控制器或边缘计算场景中。

\subsubsection{训练后量化}
\label{sec:ch04-ptq}
顾名思义，训练后量化在预训练之后施加于大语言模型。其目标通常是降低推理的内存需求，同时保持原模型的任务表现。朴素量化对权重几乎不加区分地使用低精度数据类型，已被证明对较小的语言模型有效。但大语言模型的参数量超过 67 亿时，观测到任务表现骤降，见图~\ref{fig:ch04-outlier}（\refcite{ch04-dettmers-2022}{Dettmers et al., 2022}）。这一现象与离群特征（outlier feature）的出现有关。它们表现为网络隐藏激活中的较大数值。\refcite{ch04-dettmers-2022}{Dettmers 等人（2022）}最早在大语言模型量化的背景下描述了这一现象。

要保留这些重要离群特征的表示精度，同时满足推理预算，是一个难点。\refcite{ch04-dettmers-2022}{Dettmers 等人（2022）}提出 \texttt{LLM.int8()}，对目标特征中的 99.9\% 逐向量做 INT8 量化，将离群特征单独分离，在矩阵乘法中保持 16 位精度。这增加了量化的复杂性，但这种有选择的混合精度方案使 BLOOM-176B 的推理内存需求减半，并能有效保留原大语言模型的任务表现，如图~\ref{fig:ch04-outlier} 所列的几项基准任务所示。

SqueezeLLM 采用另一种方法，以特征的二阶信息搜索最合适的位精度，从而保留离群特征和其他对精度变化敏感的特征。该方法逐层施加量化，精度可低至 3 位。与原大语言模型相比，推理速度最高可达 2.3 倍，任务表现损失很小（\refcite{ch04-kim}{Kim et al., 2023}）。

ZeroQuant 采用粒度更细的量化方式，分别为权重与激活值采用不同的量化方案，并使用一种知识蒸馏机制抵消任务表现的下降。它在准确率损失很小的情况下，最高可将推理效率提高到原来的 5 倍（\refcite{ch04-yao-2022}{Yao et al., 2022}）。

除了上述方法，GPTQ 也是一种常用的训练后量化方案。和前面的方法一样，GPTQ 用特征的二阶信息搜索量化的最合适位宽。由于有选择地量化权重，支持 4 位、3 位和 2 位的极低位宽，GPTQ 可以让具有 1760 亿参数的 BLOOM 模型在单张 NVIDIA A100 上运行，推理效率最高可达原来的 4.5 倍。\refcite{ch04-liu}{Liu 等人（2023）}则用知识蒸馏技术，尝试改善 3 位精度下的量化效果。

\sourcefigure{fig04_11.pdf}{三种数值存储精度或量化方案用于不同参数规模的语言模型，在多项自然语言推断任务中的表现。参数量从 27 亿增加到 67 亿时，表现大幅下降，这归因于离群特征的出现。\texttt{LLM.int8()} 对大语言模型权重做 8 位量化，同时保留这些特征，因此，即便存储参数权重所用精度减半，其表现仍接近完整的 16 位语言模型。}{fig:ch04-outlier}

\subsubsection{量化感知训练}
\label{sec:ch04-qat}
如第~\ref{sec:ch04-ptq}~节所述，训练后量化在学习过程中并不明确最小化量化操作引入的损失。不过，\texttt{LLM.int8()}、GPTQ 和 SqueezeLLM 等方法以较细的粒度施加训练后量化，在一定程度上也利用了学习过程中的信息。大语言模型预训练采用训练后量化的一个主要原因，是避免训练循环每次迭代时量化大量参数，显著增加计算开销。因此，许多研究试图结合量化带来的推理效率和 PEFT 带来的训练效率。与全参数微调相比，这能减少在训练时加入量化所产生的计算开销。

正如训练后量化使大语言模型能在更易获得的硬件上推理，量化感知训练也可减少微调开销，使之能在更易获得的硬件上进行（\refcite{ch04-dettmers-2023}{Dettmers et al., 2023}）。下面重点介绍原书认为有前景的基于 PEFT 的量化感知训练方法，考察它们能在多大程度上减少微调开销，以及与未经量化的 PEFT 相比，能在多大程度上保留任务表现。%
\footnote{译注：原书这里说要介绍三种方法，本小节实际列出 QLoRA 与 LoftQ 两种。}

\paragraph*{QLoRA}
\refcite{ch04-wortsman}{Wortsman 等人（2023）}提出了把推理时量化所带来的一些效率收益引入训练的思路和建议。在此基础上，QLoRA（\refcite{ch04-dettmers-2023}{Dettmers et al., 2023}）成为大语言模型中采用较广的量化感知训练方法之一。概括地说，QLoRA 先对给定大语言模型采用一种新的 4 位量化，然后在微调中冻结其参数。这项工作引入了一种新数据类型 NF4，即 4 位 NormalFloat。对于服从正态分布的张量数据，NF4 的量化精度优于 4 位整数或 4 位浮点数。量化后，LoRA 权重更新的梯度通过冻结的 4 位量化大语言模型反向传播，因此量化引入的误差会纳入微调过程。

QLoRA 除了采用上述 NF4 数据类型，还使用二次量化（double quantization）进一步减少量化常数引入的内存开销，并采用分页优化器（paged optimizer）。这些办法大幅提高了微调时的计算效率。\refcite{ch04-dettmers-2023}{Dettmers 等人（2023）}结合这三项方法，使用 LoRA 和 Alpaca、FLAN v2 数据集，对具有 650 亿参数的 Llama 大语言模型进行指令微调。他们只在一张 GPU 上微调 24 小时，所得模型即达到 ChatGPT 在相应评估中表现的 99.3\%。实际微调时，16 位精度的全参数微调需要超过 780~GB 的 GPU 显存，QLoRA 所需显存不足 48~GB，同时仍可保持接近当时最佳水平的任务表现。

\paragraph*{LoftQ}
\refcite{ch04-li-2023a}{Li 等人（2023a）}指出，预训练语言模型的量化会影响 LoRA 的微调结果，在极低位宽下尤为明显。为了减轻低位宽量化造成的精度偏差，他们提出量化感知训练技术 LoftQ，试图缩小大语言模型的原始权重矩阵与量化、低秩权重近似共同得到的权重矩阵之间的差距。

这一优化问题表述为 Frobenius 范数的最小化。
\begin{equation}
\min_{Q,A,B}\lVert W-Q-AB^{\mathsf T}\rVert_{f'},
\label{eq:ch04-loftq}
\end{equation}
其中，$\lVert\cdot\rVert_{f'}$ 表示 Frobenius 范数，$W$ 表示原始参数权重，$Q$ 表示量化后的权重，$AB^{\mathsf T}$ 表示 LoRA 所用的低秩项。

这样表述微调问题，可以为 LoRA 矩阵 $A$ 和 $B$ 找到更合适的量化初始化，也可以较好地近似原大语言模型的参数权重 $W$。为此，同时优化量化目标和 LoRA 微调目标。前者主要为了减少权重矩阵运算所需内存，后者主要为了在 $A$ 和 $B$ 的低秩约束下使任务表现尽可能好。

LoftQ 交替迭代，求解这个联合损失的最小值。先给定当前对 $A$ 和 $B$ 的估计，寻找使本步 Frobenius 范数最小的 $Q$。再给定新估计的 $Q$，用奇异值分解为 $A$ 和 $B$ 寻找量化权重 $Q$ 与原权重 $W$ 之间的残差 $Q-W$ 的低秩近似。通过交替估计量化值和执行量化感知的奇异值分解（singular value decomposition，SVD），LoftQ 在两者之间找到较好的折中，让它们共同改善微调结果。交替联合优化完成后，冻结最优的 $Q$，即可进行常规的 LoRA 微调。

LoftQ 在量化误差与 LoRA 低秩表示引入的误差之间取得了折中，而 QLoRA 并未明确最小化微调中的量化误差。由于量化相对原大语言模型引入精度偏差，QLoRA 的泛化表现不及 LoftQ。\refcite{ch04-li-2023a}{Li 等人（2023a）}测试的各项基准均显示 LoftQ 优于 QLoRA。
